% !TeX root = ../../Book1.tex
% 纲要标题：### 16.4 本原元素定理
% 纲要位置：计划纲要.md，第 576 行。

\section{本原元素定理}
\label{b1:sec:1604}


% 纲要内容（仅作写作计划，尚非教材正文）：
%
% * 先证明域的有限乘法子群为循环群；第17.2节复用此引理，不以本原元素定理反向证明循环性
% * 完整陈述本原元素定理，并连续证明无限与有限两种基域情形；解释区分全部嵌入的目标与坏参数的来源
% * 证明中的有限性与可分性
% * 非单生成有限扩张的明确反例；由指数余数对证明双变量单项式基
%
% ---
%

一个有限扩张总能由有限多个元素生成，能否进一步只用一个？可分性给出肯定答案。无限基域允许从线性组合中避开有限多个坏参数，从而把不同嵌入区分开；有限基域没有这样的参数选择余地，却能利用扩张域乘法群的循环性。两种办法分别把嵌入计数和乘法生成转化为域的单生成性。

\subsection{有限乘法子群的循环性}
\label{b1:subsec:1604:04}

\begin{lemma}\label{b1:lem:finite-multiplicative-cyclic}
任意域 \(F\) 的乘法群 \(F^\times\) 的有限子群都是循环群。
\end{lemma}
\begin{proof}
设有限子群为 \(G\)。平凡群可直接取生成元 \(1\)，以下设 \(G\ne\{1\}\)，令 \(m\) 为所有元素阶的最小公倍数。每个元素都是 \(x^m-1\) 的根，所以根数上界给出 \(\lvert G\rvert\le m\)。因此只要在 \(G\) 中构造一个阶为 \(m\) 的元素，它的 \(m\) 个幂就会占满整个群。

对每个素数幂 \(\ell^e\) 恰好整除 \(m\)，至少有一个元素 \(g_\ell\) 的阶中含有完整的 \(\ell^e\) 因子：否则所有元素阶中 \(\ell\) 的指数都小于 \(e\)，其最小公倍数也不可能含有 \(\ell^e\)。设该阶为 \(\ell^e d\)，其中 \(\ell\nmid d\)，则 \(h_\ell=g_\ell^d\) 的阶为 \(\ell^e\)：\(h_\ell^r=1\) 当且仅当 \(\ell^e d\mid dr\)，即 \(\ell^e\mid r\)。

\(G\) 交换，且各 \(h_\ell\) 的阶两两互素，所以其乘积 \(h\) 的阶为 \(m\)。为核对这一点，若两个交换元素的阶为互素的 \(a,b\)，且 \((xy)^r=1\)，则 \(x^r=y^{-r}\) 位于 \(\langle x\rangle\cap\langle y\rangle\)。交中元素的阶同时整除 \(a,b\)，故只有单位，从而 \(a\mid r,b\mid r\)，即 \(ab\mid r\)；反向 \((xy)^{ab}=1\) 直接成立。有限次应用此结论便得 \(h\) 的阶为 \(m\)。

由\cref{b1:prop:factor-root}得到的上界 \(\lvert G\rvert\le m\)，与 \(h\) 的 \(m\) 个不同幂给出的下界相等。因此 \(G=\langle h\rangle\)。
\end{proof}

\subsection{有限可分扩张的单生成性}
\label{b1:subsec:1604:01}

\begin{theorem}[本原元素]\label{b1:thm:primitive-element}
每个有限可分扩张 \(L/K\) 都是单扩张。如果 \(\theta\in L\) 满足 \(L=K(\theta)\)，则称 \(\theta\) 为该扩张的\term[benyuan-yuansu]{本原元素}[primitive element]。
\end{theorem}
\begin{proof}
先考虑无限基域 \(K\)。固定包含 \(L\) 的代数闭包 \(\Omega\)，由\cref{b1:thm:separable-embeddings}，\(L\) 有恰好 \(n=[L:K]\) 个 \(K\) 嵌入 \(\sigma_1,\ldots,\sigma_n\)。我们要找到一个元素，使这些嵌入在它上面取值两两不同；这样一个元素的最小多项式就必须容纳 \(n\) 个不同根，其次数才能达到整个扩张的次数。

先设 \(L=K(a,b)\)，考察 \(a+cb\)、\(c\in K\)。对于 \(i\neq j\)，若 \(\sigma_i(b)=\sigma_j(b)\)，必有 \(\sigma_i(a)\neq\sigma_j(a)\)，否则它们在两个生成元上相同，因而为同一个映射。此时
\[
 \sigma_i(a+cb)-\sigma_j(a+cb)=\sigma_i(a)-\sigma_j(a)\ne0,
\]
所以每个 \(c\in K\) 都把这一对嵌入区分开。若 \(\sigma_i(b)\neq\sigma_j(b)\)，等式
\[
\sigma_i(a)+c\cdot\sigma_i(b)=\sigma_j(a)+c\cdot\sigma_j(b)
\]
只可能在
\[
 c=\frac{\sigma_j(a)-\sigma_i(a)}{\sigma_i(b)-\sigma_j(b)}
\]
时成立。这个值若不在 \(K\) 中，就不排除任何系数；若在 \(K\) 中，也只排除一个。嵌入对只有有限多个，坏系数也只有有限多个。这里才用到 \(K\) 无限：这些值不能耗尽 \(K\)，所以可以同时避开它们。

令 \(\theta=a+cb\)。每个 \(\sigma_i\) 固定最小多项式的系数，所以
\[
 m_{\theta,K}(\sigma_i(\theta))
 =\sigma_i\bigl(m_{\theta,K}(\theta)\bigr)=0.
\]
这些值两两不同，故 \(m_{\theta,K}\) 至少有 \(n\) 个不同根。由次数与根数的关系及塔式公式，
\[
n\leq\bigl[K(\theta):K\bigr]\leq[L:K]=n.
\]
所以 \(\bigl[L:K(\theta)\bigr]=1\)，即 \(L=K(\theta)\)。一般情形取有限生成集 \(a_1,\ldots,a_r\)。合并前两个后得到 \(\theta_2\)，满足 \(K(\theta_2)=K(a_1,a_2)\)。因为 \(\theta_2\in L\) 且 \(L/K\) 可分，\(\theta_2\) 仍在 \(K\) 上可分；于是 \(K(\theta_2,a_3)/K\) 仍是有限可分扩张，可以再次应用二生成元情形。这样逐次合并，便得到生成整个 \(L\) 的一个元素。

若基域有限，有限多个坏值有可能占满整个基域，因而不能再靠上面的避让选出系数。记 \(\lvert K\rvert=q\)、\([L:K]=n\)。选一组 \(K\) 基，每个元素由 \(n\) 个 \(K\) 坐标唯一确定，所以 \(\lvert L\rvert=q^n\)，是有限域。\cref{b1:lem:finite-multiplicative-cyclic}说明 \(L^\times=\langle\theta\rangle\)。域 \(K(\theta)\) 包含 \(\theta\) 的全部整数幂，因而包含 \(L\) 的每个非零元素；它也包含零，故等于 \(L\)。这实际上证明了有限域的任意有限扩张都单生成，论证本身不需要可分性。有限域的完美性已由\cref{b1:prop:perfect-frobenius}得到，所以这些扩张也都可分。
\end{proof}
这一结果称为\term[benyuan-yuansu-dingli]{本原元素定理}[primitive element theorem]。有限域乘法群的生成元一定是域扩张的本原元素，反之未必；前者必须用乘法的幂得到所有非零元素，后者允许加法、减法和除法。

\ProseParagraphBreak
例如 \(X^2+1\) 在 \(\mathbb F_3\) 中无根，故 \(F=\mathbb F_3[X]/(X^2+1)\) 是九元素域。令 \(a\) 为 \(X\) 的剩余类，则 \(F=\mathbb F_3(a)\)，但 \(a^2=-1\ne1\)、\(a^4=1\)，所以 \(a\) 的乘法阶为四，不能生成八阶的 \(F^\times\)。

\begin{corollary}\label{b1:prop:primitive-element-infinite}
设 \(K\) 为无限域，\(L/K\) 为有限可分扩张，\(a,b\in L\) 且 \(L=K(a,b)\)。记 \(n=[L:K]\)。则使 \(L\ne K(a+cb)\) 的系数 \(c\in K\) 至多有 \(\binom{n}{2}\) 个。
\end{corollary}
\begin{proof}
取包含 \(L\) 的 \(K\) 的代数闭包 \(\Omega\)，沿用\cref{b1:thm:primitive-element}对无限基域的证明。对 \(n\) 个不同的 \(K\) 嵌入 \(L\rightarrow\Omega\) 中的每一无序对，等式 \(\sigma_i(a+cb)=\sigma_j(a+cb)\) 或无解，或在 \(K\) 中至多有一个解。这样的嵌入对共 \(\binom{n}{2}\) 个；避开这些解时，全部嵌入在 \(a+cb\) 上取值不同，前一定理的次数论证便给出 \(L=K(a+cb)\)。
\end{proof}

例如 \(\theta=\sqrt2+\sqrt3\) 生成 \(\mathbb Q(\sqrt2,\sqrt3)\)。这不只由存在性定理得出，还可直接恢复两个根：\(\theta^{-1}=\sqrt3-\sqrt2\)，从而
\[
\sqrt2=\frac{\theta-\theta^{-1}}2,\qquad
\sqrt3=\frac{\theta+\theta^{-1}}2.
\]

\subsection{有限性与可分性的作用}
\label{b1:subsec:1604:02}

在无限基域的证明中，有限性保证嵌入数有限，因而只有有限多个嵌入对和坏参数；可分性则保证嵌入数恰为 \([L:K]\)。所以，把这些嵌入在一个元素上的值全部分开，才能迫使该元素的最小多项式次数达到 \([L:K]\)。若嵌入数小于次数，这一办法只给出较弱的次数下界，不能推出该元素生成整个扩张。

\ProseParagraphBreak
无限次代数扩张不可能单生成：若 \(L/K\) 是这样的扩张且 \(L=K(a)\)，则 \(a\) 在 \(K\) 上代数，由\cref{b1:prop:minimal-polynomial-degree}知 \([L:K]=\deg m_{a,K}<\infty\)，矛盾。因此，即使扩张可分，也不能省去本原元素定理的有限性假设。例如，\cref{b1:prop:infinite-algebraic-root-tower}中的 \(\bigcup_{n\geq1}\mathbb Q(2^{1/2^n})/\mathbb Q\) 是无限次代数扩张，并因特征为零而可分。

\ProseParagraphBreak
另一方面，可分性是充分条件，并非单生成的必要条件。\(\mathbb F_p(t^{1/p})/\mathbb F_p(t)\) 本来就是单扩张，却纯不可分。因此有限可分推出单生成，单生成却未必可分；而下面的例子表明，有限纯不可分扩张也未必单生成。

\subsection{非单生成有限扩张}
\label{b1:subsec:1604:03}

\begin{proposition}\label{b1:prop:non-simple-finite-extension}
设 \(p\) 为素数，\(u,v\) 在 \(\mathbb F_p\) 上代数独立，并令
\[
L=\mathbb F_p(u,v),\qquad K=\mathbb F_p(u^p,v^p).
\]
则 \([L:K]=p^2\)，但对每个 \(a\in L\) 都有 \(\bigl[K(a):K\bigr]\leq p\)。特别地，这个有限扩张不是单扩张。
\end{proposition}
\begin{proof}
因为 \(u^p,v^p\in K\)，高次幂的指数都能模 \(p\) 降低，所以自然考察 \(p\times p\) 个单项式
\[
 \{\,u^iv^j:0\le i,j<p\,\}.
\]
\(u,v\) 分别满足 \(x^p-u^p\)、\(x^p-v^p\)，所以总次数至多 \(p^2\)。要证明这些单项式是一组基，只需证明它们在 \(K\) 上线性无关。若有关系，清除系数中 \(u^p,v^p\) 的有理函数分母，得到
\[
 \sum_{0\le i,j<p}P_{ij}(u^p,v^p)u^iv^j=0,
 \qquad P_{ij}\in\mathbb F_p[U,V].
\]
展开第 \((i,j)\) 项后，每个单项式的指数对形如 \((pm+i,pn+j)\)。不同 \((i,j)\) 给出不同的模 \(p\) 余数对，所以这些项所含的单项式集合互不相交。\(u,v\) 在 \(\mathbb F_p\) 上代数独立，迫使各 \(P_{ij}\) 的全部系数为零，进而迫使原关系的系数全为零。于是扩张次数恰为 \(p^2\)，这些单项式组成一组 \(K\) 基。

对任意 \(a=A(u,v)/B(u,v)\)，其中 \(A,B\in\mathbb F_p[U,V]\)、\(B(u,v)\ne0\)，特征 \(p\) 使取幂时的交叉项消失，且系数满足 \(c^p=c\)，所以
\[
 a^p=\frac{A(u,v)^p}{B(u,v)^p}
     =\frac{A(u^p,v^p)}{B(u^p,v^p)}\in K.
\]
分母仍非零，因为域中非零元素的幂非零。因此 \(a\) 被次数为 \(p\) 的多项式 \(x^p-a^p\in K[x]\) 零化，给出 \(\bigl[K(a):K\bigr]\leq p<p^2\)。因而任何单个 \(a\) 都不能生成 \(L\)。
\end{proof}
这里所说的两个 \(p\) 次根方向，是指 \(u\) 与 \(v\) 各带来一个次数为 \(p\) 的步骤，总次数为 \(p^2\)，而每个元素单独只能带来至多 \(p\) 的次数。它不是说 \(u,v\) 在 \(K\) 上代数独立；二者在 \(K\) 上都已代数。
